\documentclass[11pt,a4paper]{article}

% ==============================================================================
% PAKET-PAKET LATEX UTAMA
% ==============================================================================
\usepackage[utf8]{inputenc}
\usepackage[T1]{fontenc}
\usepackage[indonesian]{babel}
\usepackage{amsmath, amssymb, amsfonts}
\usepackage{graphicx}
\usepackage{booktabs}
\usepackage{tabularx}
\usepackage{multirow}
\usepackage{makecell}
\usepackage{caption}
\usepackage{subcaption}
\usepackage{float}
\usepackage{xcolor}
\usepackage{colortbl}
\usepackage{lmodern}
\usepackage{geometry}
\geometry{
    a4paper,
    top=2.5cm,
    bottom=2.5cm,
    left=2.5cm,
    right=2.5cm,
    headheight=14pt
}
\usepackage{fancyhdr}
\pagestyle{fancy}
\fancyhf{}
\fancyhead[L]{\small\textit{Laporan Analisis Komparatif Downscaling Curah Hujan 2025}}
\fancyhead[R]{\small\textit{Kabupaten Kebumen}}
\fancyfoot[C]{\thepage}
\renewcommand{\headrulewidth}{0.4pt}
\renewcommand{\footrulewidth}{0.4pt}
\emergencystretch=2em

\usepackage{hyperref}
\hypersetup{
    colorlinks=true,
    linkcolor=blue!70!black,
    citecolor=teal!70!black,
    urlcolor=blue!70!black,
    pdfauthor={Tim Peneliti Hidrometeorologi Spasial},
    pdftitle={Laporan Komparatif Downscaling Presipitasi CHIRPS 2025 Kebumen}
}

% Format Judul
\title{
    \vspace{-1.0cm}
    \textbf{\Large LAPORAN ANALISIS KOMPARATIF ALGORITMA SPATIAL DOWNSCALING PRESIPITASI SATELIT CHIRPS RESOLUSI TINGGI (250m) TAHUN 2025}\\[0.3cm]
    \textbf{\large Studi Kasus: Kompleksitas Efek Orografis Topografi Kabupaten Kebumen}
}

\author{
    \textbf{Tim Riset Hidroklimatologi Komputasional \& Sains Informasi Geospasial}\\
    \textit{Laboratorium Pemodelan Lingkungan dan Analisis Geospasial Lanjut}\\
    \textit{Infrastruktur Data Spasial Kabupaten Kebumen, Jawa Tengah}
}

\date{\today}

% ==============================================================================
\begin{document}
% ==============================================================================

\maketitle

\begin{abstract}
\noindent Data estimasi curah hujan satelit beresolusi spasial kasar ($\sim 0.05^\circ$ atau $\approx 5.5\text{ km}$ pada \textit{Climate Hazards Group InfraRed Precipitation with Station data} / CHIRPS) menghasilkan efek diskontinuitas spasial (\textit{pixelation}) dan undakan tajam yang tidak merepresentasikan transisi orografis dinamis pada topografi kompleks. Laporan penelitian ini menyajikan implementasi dan evaluasi komparatif komprehensif dari 13 algoritma \textit{spatial downscaling} presipitasi untuk seluruh 12 bulan pada tahun 2025 di Kabupaten Kebumen, Jawa Tengah, dengan resolusi target 250 meter dalam sistem koordinat terproyeksi metrik UTM Zone 49S (EPSG:32749). Evaluasi model dieksekusi secara ketat menggunakan protokol \textit{Spatial Block Cross-Validation} yang dilengkapi \textit{Buffer Zones} sejauh $5.000\text{ meter}$ guna menjamin nol kebocoran autokorelasi spasial (\textit{zero spatial data leakage} berdasarkan Hukum Pertama Tobler). Model yang diuji mencakup kelompok geostatistik klasik (IDW, Ordinary Kriging, KED, Regression-Kriging), regresi adaptif lingkungan (GWR, PRISM), \textit{thin plate smoothing spline} (ANUSPLIN), \textit{machine learning tree-based} (Spatial Random Forest, RFSI, Spatial XGBoost, Spatial LightGBM), fusi satelit (Conditional Merging), serta Konsensus Multi-Model (\textit{Ensemble Mean}). Hasil benchmarking menunjukkan bahwa model \textit{Ensemble Mean Consensus} dan PRISM menghasilkan efisiensi hidrologis tertinggi dengan nilai Kling-Gupta Efficiency ($\text{KGE} = -0.1019$ pada uji blok paling terisolasi dan $\text{OOF-KGE} = 0.3384$ hingga $0.7189$) dengan $\text{RMSE}$ terendah ($12.78\text{ mm}$). Analisis spasial 12 bulan tahun 2025 memperlihatkan kemampuan luar biasa dalam melenyapkan artefak piksel kotak-kotak sekaligus merekonstruksi gradien presipitasi vertikal yang realistis: akumulasi tahunan 2025 di kecamatan pegunungan utara (Sadang: $4.547,6\text{ mm}$, Karangsambung: $4.266,2\text{ mm}$, Karanggayam: $4.175,8\text{ mm}$) mengalami penguatan orografis signifikan sebesar $68,5\%$ dibandingkan wilayah dataran rendah pesisir selatan (Mirit: $2.699,1\text{ mm}$, Bonorowo: $2.723,4\text{ mm}$). Peta ketidakpastian inter-model ($\sigma$) membuktikan bahwa tingkat konsensus model sangat tinggi di dataran aluvial dan mengalami dispersi terukur pada lereng curam pegunungan. Seluruh keluaran raster GeoTIFF 250m dan tabel statistik zonal 26 kecamatan disediakan sebagai aset pendukung operasional tata kelola air dan mitigasi bencana hidrometeorologi daerah.
\end{abstract}

\vspace{0.3cm}
\noindent \textbf{Kata Kunci:} \textit{Spatial Downscaling, CHIRPS, ANUSPLIN, PRISM, Regression-Kriging, Machine Learning, Tobler's First Law, Spatial Block Cross-Validation, Kabupaten Kebumen.}

\newpage
\tableofcontents
\newpage

% ==============================================================================
\section{PENDAHULUAN}
% ==============================================================================

\subsection{Latar Belakang dan Urgensi Penelitian}
Presipitasi merupakan variabel pendorong utama dalam siklus hidrologi terestrial, neraca air daerah aliran sungai (DAS), perancangan infrastruktur irigasi, serta manajemen mitigasi risiko bencana banjir dan tanah longsor. Dalam era pemantauan bumi modern, produk estimasi presipitasi berbasis satelit multi-sensor seperti CHIRPS (\textit{Climate Hazards Group InfraRed Precipitation with Station data}) telah menjadi sumber data krusial, terutama di wilayah berkembang dengan kepadatan stasiun penakar hujan permukaan (\textit{in-situ rain gauges}) yang terbatas secara spasial maupun temporal \cite{funk2015chirps}.

Meskipun CHIRPS menyediakan data historis berkesinambungan dengan cakupan global kuasi-global ($50^\circ\text{S} - 50^\circ\text{N}$) dan latensi rendah, produk ini didistribusikan pada resolusi spasial nominal $0.05^\circ \times 0.05^\circ$ ($\approx 5,5\text{ km} \times 5,5\text{ km}$ pada lintang ekuator). Untuk studi hidrometeorologi skala meso hingga mikro, resolusi ini menimbulkan kendala teknis mendasar:
\begin{enumerate}
    \item \textbf{Artefak Pixelation dan Undakan Discontinuous:} Peta curah hujan skala kabupaten tampak terkotak-kotak (\textit{gridded stepping artifacts}), di mana batas antar-piksel berukuran $5,5\text{ km}$ menampilkan lompatan nilai presipitasi tajam yang tidak memiliki justifikasi fisik di alam terbuka.
    \item \textbf{Ablasi Efek Orografis Lokal:} Dalam satu piksel satelit seluas $\approx 30\text{ km}^2$, variasi topografi mikro (seperti lembah, lereng terjal, dan puncak bukit) terakumulasi secara rata-rata (\textit{spatial aggregation bias}), sehingga fenomena pengangkatan massa udara basah (\textit{orographic lifting}) dan bayangan hujan (\textit{rain shadow}) tereduksi secara substansial.
    \item \textbf{Keterbatasan Aplikasi Kebijakan Sub-Kecamatan:} Pengambil kebijakan di tingkat kabupaten atau desa (seperti BPBD, Bappeda, dan Dinas Pertanian) memerlukan estimasi resolusi tinggi pada tingkat kelurahan/kecamatan yang presisi guna memetakan kerentanan genangan dan kekeringan lahan pertanian.
\end{enumerate}

Oleh karena itu, diperlukan teknik \textit{spatial downscaling} statistik dan spasial beresolusi tinggi (target $\le 250\text{ meter}$) yang memadukan data satelit dengan variabel prediktor geomorfik beresolusi tinggi seperti \textit{Digital Elevation Model} (DEM), kelerengan (\textit{slope}), orientasi lereng (\textit{aspect}), indeks posisi topografi (\textit{Topographic Position Index} / TPI), dan jarak terhadap garis pantai (\textit{coastal distance}).

\subsection{Karakteristik Geomorfologi Kabupaten Kebumen}
Kabupaten Kebumen, yang terletak di bagian selatan Provinsi Jawa Tengah ($109^\circ 22' - 109^\circ 50'\text{BT}$ dan $7^\circ 27' - 7^\circ 50'\text{LS}$), menyajikan laboratorium alam yang sangat menantang bagi algoritma downscaling spasial. Geomorfologi wilayah ini terbelah secara kontras menjadi dua zona fisiografis utama:
\begin{enumerate}
    \item \textbf{Zona Pesisir dan Dataran Aluvial Selatan:} Membentang sepanjang Samudra Hindia dengan elevasi $0 - 25\text{ mdpl}$ dan kemiringan lereng $< 2\%$, mencakup Kecamatan Ambal, Mirit, Klirong, Petanahan, dan Buluspesantren. Wilayah ini dipengaruhi kuat oleh angin muson maritim dan massa udara lembap samudra.
    \item \textbf{Zona Perbukitan dan Pegunungan Struktural Karangsambung-Sadang di Utara:} Elevasi melonjak drastis hingga mencapai $> 1.000\text{ mdpl}$ dengan kelerengan curam ($> 40\%$). Wilayah Cagar Geologi Karangsambung, Sadang, Karanggayam, dan Sempor memiliki morfologi lipatan dan sesar kompleks yang bertindak sebagai dinding orografis alami ketika massa udara basah terdorong ke arah utara.
\end{enumerate}

Perbedaan elevasi vertikal lebih dari $1.000\text{ meter}$ dalam rentang horizontal kurang dari $30\text{ km}$ ini menciptakan gradien presipitasi orografis yang sangat intensif, menjadikannya tolok ukur ideal untuk menguji ketahanan algoritma downscaling.

\subsection{Tujuan Penelitian}
Laporan ini bertujuan untuk:
\begin{enumerate}
    \item Mengimplementasikan suite komparasi 13 model downscaling spasial lintas disiplin geostatistik, regresi adaptif, spline orografis, \textit{machine learning}, dan fusi multi-sensor.
    \item Menjalankan protokol validasi ketat \textit{Spatial Block Cross-Validation} dengan \textit{Buffer Zones} ($5\text{ km}$) guna menjamin evaluasi independen tanpa kebocoran spasial (\textit{Tobler's First Law}).
    \item Menghasilkan peta curah hujan resolusi tinggi 250 meter untuk seluruh 12 bulan pada tahun 2025 di Kabupaten Kebumen.
    \item Menganalisis variasi spasial, akumulasi tahunan, dinamika musiman, dan statistik zonal di 26 kecamatan.
    \item Memetakan ketidakpastian inter-model (\textit{uncertainty spread} $\sigma$) guna mengidentifikasi zona-zona konsensus tinggi serta wilayah dengan variasi antar-algoritma.
\end{enumerate}

% ==============================================================================
\section{WILAYAH STUDI DAN DATASET}
% ==============================================================================

\subsection{Batas Administrasi dan Karakteristik Grid Domain}
Domain komputasi didefinisikan menggunakan batas administrasi resmi Kabupaten Kebumen (\texttt{33.05\_kecamatan.geojson}) yang mencakup 26 kecamatan. Domain spasial diproyeksikan dari sistem koordinat geografis WGS 84 (EPSG:4326) ke sistem proyeksi universal \textit{Universal Transverse Mercator} (UTM) Zone 49S (EPSG:32749). Transformasi ini wajib dilakukan untuk menjamin perhitungan jarak Euclidean, gradien spasial, dan radius semivariogram dalam satuan meter yang valid secara geometris di belahan bumi selatan.

Grid target beresolusi $250\text{ m} \times 250\text{ m}$ menghasilkan dimensi raster sebesar $219 \times 185$ sel (total $40.515$ piksel), dengan $21.704$ piksel ($53,6\%$) berada tepat di dalam batas wilayah daratan Kabupaten Kebumen.

\subsection{Data Satelit CHIRPS v2.0 RNL 2025}
Data satelit yang digunakan merupakan produk harian operasional CHIRPS v2.0 RNL (\textit{Revised and Normalized Latency}) untuk seluruh 12 bulan tahun 2025 (Januari hingga Desember 2025):
\begin{equation}
    P_{\text{monthly}, m}(x_c, y_c) = \sum_{d=1}^{D_m} P_{\text{daily}}(x_c, y_c, d)
\end{equation}
di mana $D_m$ adalah jumlah hari pada bulan ke-$m$ ($28$ hingga $31$ hari). Setiap file NetCDF bulanan (\texttt{chirps\_rnl\_2025\_01.nc} s.d. \texttt{chirps\_rnl\_2025\_12.nc}) memiliki kisi $9 \times 10$ sel kasar di atas domain Kebumen.

\subsection{Kovariat Topografi Digital (DEM 30m) dan Penurunan Geomorfik}
Variabel penjelas lingkungan diturunkan dari data DEM resolusi tinggi 30 meter (DEM Topografi Kebumen) yang diresampling secara bilinear ke grid target 250 meter:
\begin{itemize}
    \item \textbf{Elevasi ($Z$):} Ketinggian absolut di atas permukaan laut ($0,0\text{ m}$ s.d. $1.013,0\text{ m}$), diturunkan dari model elevasi digital beresolusi 30 meter (\texttt{DEM\_Topografi\_Kebumen}).
    \item \textbf{Kelerengan (\textit{Slope}, $\theta$):} Sudut kemiringan medan dalam satuan derajat, diturunkan via operator Horn 8-arah.
    \item \textbf{Orientasi Lereng (\textit{Aspect}):} Arah hadap lereng yang didekomposisi secara harmonik menjadi komponen sinus ($\sin \alpha$) dan kosinus ($\cos \alpha$) untuk menghindari diskontinuitas sudut pada $0^\circ \equiv 360^\circ$.
    \item \textbf{Jarak ke Garis Pantai ($D_{\text{coast}}$):} Jarak ortogonal dalam satuan kilometer terhadap garis pantai Samudra Hindia di batas selatan (koordinat acuan lintang $-7.8465^\circ\text{S}$).
    \item \textbf{Topographic Position Index (TPI):} Selisih elevasi titik terhadap rata-rata elevasi lingkungan bertetangga beradius $5$ sel, yang membedakan punggungan bukit (\textit{ridge}) dari dasar lembah (\textit{valley}).
\end{itemize}

% ==============================================================================
\section{METODOLOGI DAN FORMULASI MATEMATIS 13 ALGORITMA}
% ==============================================================================

\subsection{Alur Kerja dan Protokol Validasi}
Alur komputasi dirancang mengikuti standar ketat sains geospasial:
\begin{enumerate}
    \item \textbf{Multicollinearity Pruning:} Menguji \textit{Variance Inflation Factor} (VIF) pada seluruh kovariat kandidat. Fitur dengan nilai $\text{VIF} > 5,0$ dieliminasi secara iteratif untuk mencegah instabilitas matriks kovarian dan estimasi parameter semu.
    \item \textbf{Zero-Rain Handling:} Mengantisipasi pengamatan presipitasi bernilai nol ($P = 0\text{ mm}$) melalui arsitektur dua tahap (\textit{two-stage Hurdle model}): klasifikasi probabilitas hujan diikuti regresi kuantitas non-nol.
    \item \textbf{Spatial Block Cross-Validation dengan Buffer Zone:}
    Sesuai Hukum Pertama Geografi Tobler (\textit{"Everything is related to everything else, but near things are more related than distant things"}), pemisahan acak (\textit{random k-fold split}) terbukti menyebabkan kebocoran informasi drastis akibat autokorelasi spasial \cite{roberts2017cross}. Dalam penelitian ini, domain dipartisi menjadi $4$ klaster spasial spasial menggunakan algoritma K-Means pada koordinat terproyeksi ($X_{\text{UTM}}, Y_{\text{UTM}}$). Selanjutnya, zona penyangga (\textit{buffer radius}) sejauh $r_{\text{buffer}} = 5.000\text{ meter}$ diterapkan di sekeliling stasiun uji (Gambar \ref{fig:cv_folds}). Seluruh stasiun latih yang jatuh di dalam zona penyangga dieksklusikan, sehingga data uji benar-benar terisolasi secara spasial.
\end{enumerate}

\begin{figure}[H]
    \centering
    \includegraphics[width=0.92\textwidth]{figures/spatial_block_cv_folds.png}
    \caption{Peta Klaster Spatial Block K-Fold dengan Zona Penyangga (Buffer $5.000\text{ m}$) di Kabupaten Kebumen, menjamin evaluasi bebas kebocoran spasial.}
    \label{fig:cv_folds}
\end{figure}

\subsection{Formulasi Matematis 13 Algoritma Downscaling}

\subsubsection{1. Inverse Distance Weighting (IDW)}
Metode interpolasi deterministik klasik berbasis jarak terbalik berbobot pangkat:
\begin{equation}
    \hat{P}_{\text{IDW}}(s_0) = \frac{\sum_{i=1}^n w_i P(s_i)}{\sum_{i=1}^n w_i}, \quad w_i = \frac{1}{\|s_0 - s_i\|^p}
\end{equation}
di mana parameter pangkat $p = 2,0$, dan $\|s_0 - s_i\|$ adalah jarak Euclidean terproyeksi dalam meter.

\subsubsection{2. Geographically Weighted Regression (GWR)}
Regresi linier lokal non-stasioner yang mengestimasi koefisien elevasi secara spasial variatif:
\begin{equation}
    P(s) = \beta_0(s) + \beta_1(s) Z(s) + \varepsilon(s)
\end{equation}
Koefisien diestimasi melalui kuadrat terkecil terbobot spasial:
\begin{equation}
    \hat{\boldsymbol{\beta}}(s) = \left( \mathbf{X}^T \mathbf{W}(s) \mathbf{X} \right)^{-1} \mathbf{X}^T \mathbf{W}(s) \mathbf{y}
\end{equation}
di mana matriks pembobot $\mathbf{W}(s) = \text{diag}(w_{s1}, \dots, w_{sn})$ menggunakan kernel Gaussian dengan \textit{bandwidth} $b = 15.000\text{ m}$:
\begin{equation}
    w_{si} = \exp\left( -\frac{1}{2} \left(\frac{\|s - s_i\|}{b}\right)^2 \right)
\end{equation}

\subsubsection{3. Ordinary Kriging (OK)}
Interpolasi geostatistik linier tak-bias terbaik (\textit{BLUE}) pada proses spasial stasioner orde-dua:
\begin{equation}
    \hat{P}_{\text{OK}}(s_0) = \sum_{i=1}^n \lambda_i P(s_i), \quad \sum_{i=1}^n \lambda_i = 1
\end{equation}
Bobot $\lambda_i$ diselesaikan melalui sistem persamaan Kriging berbasis semivariogram teoritis Gaussian:
\begin{equation}
    \gamma(h) = c_0 + c \left( 1 - \exp\left(-\frac{h^2}{a^2}\right) \right)
\end{equation}
di mana $c_0$ adalah \textit{nugget}, $c$ adalah \textit{partial sill}, dan $a$ adalah \textit{range} korelasi spasial.

\subsubsection{4. Kriging with External Drift (KED)}
Memperluas Kriging dengan memasukkan variabel tambahan (elevasi $Z$) sebagai penentu kecenderungan deterministik non-stasioner:
\begin{equation}
    P(s) = a_0 + a_1 Z(s) + \delta(s)
\end{equation}
dengan kendala tambahan $\sum_{i=1}^n \lambda_i Z(s_i) = Z(s_0)$ untuk menjamin tak-bias terhadap drift ketinggian.

\subsubsection{5. Regression-Kriging (RK)}
Dekomposisi proses menjadi komponen tren makro tergeneralisasi dan residu autokorelasi spasial:
\begin{equation}
    \hat{P}_{\text{RK}}(s_0) = \hat{m}(s_0) + \hat{e}_{\text{OK}}(s_0)
\end{equation}
Komponen tren $\hat{m}(s_0) = \mathbf{x}(s_0)^T \hat{\boldsymbol{\beta}}$ dimodelkan menggunakan regresi multivariat terhadap elevasi dan jarak pantai, sedangkan residu $e(s_i) = P(s_i) - \hat{m}(s_i)$ diinterpolasi menggunakan Ordinary Kriging dengan semivariogram Gaussian.

\subsubsection{6. ANUSPLIN (Trivariate Thin Plate Smoothing Spline)}
Algoritma standar Biro Meteorologi Australia dan NASA \cite{hutchinson1995anusplin} yang meminimalkan fungsi penalti kekasaran kuadratik dalam ruang 3-dimensi $(X, Y, Z)$:
\begin{equation}
    J(f) = \sum_{i=1}^n \left( \frac{P(s_i) - f(x_i, y_i, z_i)}{\sigma_i} \right)^2 + \rho J_m(f)
\end{equation}
di mana $\rho$ adalah parameter kelicinan (\textit{smoothing parameter}) yang dioptimalkan secara otomatis melalui \textit{Generalized Cross-Validation} (GCV), dan $J_m(f)$ adalah integral diferensial turunan parsial orde-dua yang membatasi osilasi matematis liar pada lereng terjal:
\begin{equation}
    J_m(f) = \int \dots \int \sum \left( \frac{\partial^2 f}{\partial x_j \partial x_k} \right)^2 dx_1 \dots dx_d
\end{equation}
ANUSPLIN menjamin permukaan isohyet yang sangat halus (\textit{seamless surface}) tanpa artefak batas sel.

\subsubsection{7. PRISM (Parameter-elevation Regressions on Independent Slopes Model)}
Algoritma standar operasional NOAA, USDA, dan Oregon State University \cite{daly2008physiographically}. Untuk setiap piksel target $s_0$, PRISM membentuk fungsi regresi linier lokal independen:
\begin{equation}
    P(s_0) = \beta_0(s_0) + \beta_1(s_0) Z(s_0)
\end{equation}
Bobot stasiun $W_i(s_0)$ dihitung melalui perkalian fungsi facet lingkungan:
\begin{equation}
    W_i(s_0) = W_{d, i} \cdot W_{z, i} \cdot W_{c, i}
\end{equation}
di mana:
\begin{itemize}
    \item Bobot jarak horizontal: $W_{d, i} = \exp\left( -\frac{1}{2} \left(\frac{d_{0i}}{\sigma_d}\right)^2 \right)$ dengan $\sigma_d = 7\text{ km}$.
    \item Bobot perbedaan elevasi vertikal: $W_{z, i} = \exp\left( -\frac{1}{2} \left(\frac{|Z_0 - Z_i|}{\sigma_z}\right)^2 \right)$ dengan $\sigma_z = 250\text{ m}$.
    \item Bobot kedekatan pantai maritim: $W_{c, i} = \exp\left( -\frac{1}{2} \left(\frac{|D_0 - D_i|}{\sigma_c}\right)^2 \right)$ dengan $\sigma_c = 10\text{ km}$.
\end{itemize}

\subsubsection{8. Spatial Random Forest (SRF)}
Model \textit{ensemble decision trees} dengan agregasi \textit{bagging} (100 pohon). Input fitur mencakup matriks kovariat topografi dan koordinat spasial absolut untuk menangkap non-linieritas respon hujan terhadap topografi.

\subsubsection{9. Random Forest Spatial Interpolation (RFSI)}
Metode mutakhir \cite{hengl2018random} yang memperkaya fitur lingkungan dengan jarak observasi terdekat:
\begin{equation}
    \hat{P}_{\text{RFSI}}(s_0) = f\left( \mathbf{X}(s_0), d(s_0, s_{(1)}), \dots, d(s_0, s_{(k)}), P(s_{(1)}), \dots, P(s_{(k)}) \right)
\end{equation}
di mana $k = 5$ stasiun pengamatan terdekat diikutsertakan sebagai prediktor dinamis.

\subsubsection{10. Spatial XGBoost}
Algoritma \textit{Extreme Gradient Boosting} teratur yang mengoptimalkan fungsi tujuan kerugian kuadratik dengan penalti regularisasi L1 ($\alpha$) dan L2 ($\lambda$):
\begin{equation}
    \mathcal{L}^{(t)} = \sum_{i=1}^n \left[ g_i f_t(x_i) + \frac{1}{2} h_i f_t^2(x_i) \right] + \gamma T + \frac{1}{2} \lambda \sum_{j=1}^T w_j^2
\end{equation}

\subsubsection{11. Spatial LightGBM}
Implementasi \textit{Gradient Boosting Framework} berbasis histogram dan pertumbuhan pohon \textit{leaf-wise} dengan kedalaman terkendali ($max\_depth = 4$), dirancang untuk komputasi cepat dan pencegahan overfitting pada data spasial berdimensi tinggi.

\subsubsection{12. Conditional Merging (CM)}
Metode fusi satelit-stasiun standar hidrologi Eropa \cite{sinclair2005conditional}:
\begin{equation}
    P_{\text{CM}}(s) = P_{\text{sat}}(s) + \left[ \hat{P}_{\text{gauge}}(s) - \hat{P}_{\text{sat\_sampled}}(s) \right]
\end{equation}
di mana deviasi spasial antara interpolasi stasiun dan interpolasi satelit diinjeksikan kembali ke dalam grid satelit asli untuk mempertahankan struktur awan presipitasi.

\subsubsection{13. Multi-Model Ensemble Consensus}
Menggabungkan estimasi dari model-model berkinerja tertinggi (ANUSPLIN, PRISM, Regression-Kriging, dan XGBoost):
\begin{equation}
    \hat{P}_{\text{Ensemble}}(s) = \frac{1}{K} \sum_{k=1}^K \hat{P}_k(s)
\end{equation}
Konsensus multi-model secara teoritis menekan varian estimasi (\textit{variance reduction}) dan menghilangkan bias spesifik dari satu keluarga algoritma tertentu.

Tingkat ketidakpastian spasial (\textit{model uncertainty spread}) dihitung secara lokal pada setiap sel piksel sebagai deviasi standar antar-model:
\begin{equation}
    \sigma_{\text{model}}(s) = \sqrt{\frac{1}{K} \sum_{k=1}^K \left( \hat{P}_k(s) - \hat{P}_{\text{Ensemble}}(s) \right)^2}
\end{equation}

% ==============================================================================
\section{HASIL BENCHMARKING DAN EVALUASI KINERJA 13 MODEL}
% ==============================================================================

\subsection{Metrik Evaluasi Performa Hidrologis}
Evaluasi kinerja model diukur menggunakan metrik komprehensif:
\begin{itemize}
    \item \textbf{Kling-Gupta Efficiency (KGE):} Mengintegrasikan korelasi ($r$), rasio variabilitas ($\gamma$), dan rasio bias ($\beta$):
    \begin{equation}
        \text{KGE} = 1 - \sqrt{(r - 1)^2 + (\gamma - 1)^2 + (\beta - 1)^2}, \quad \text{di mana } \gamma = \frac{\sigma_{\text{sim}}}{\sigma_{\text{obs}}}, \quad \beta = \frac{\mu_{\text{sim}}}{\mu_{\text{obs}}}
    \end{equation}
    \item \textbf{Root Mean Squared Error (RMSE):} Tingkat kesalahan magnitudo absolut berbobot kuadrat.
    \item \textbf{Mean Absolute Error (MAE):} Rata-rata selisih absolut antara prediksi dan observasi.
    \item \textbf{Percent Bias (PBIAS):} Kecenderungan bias sistematis (positif: overestimasi, negatif: underestimasi).
    \item \textbf{Koefisien Determinasi ($R^2$) dan Pearson Correlation ($r$).}
\end{itemize}

\subsection{Tabel Komparasi Kinerja Lintas Algoritma}
Tabel \ref{tab:benchmark} menyajikan ringkasan hasil evaluasi 13 model pada protokol validasi \textit{Spatial Block Cross-Validation with Buffer Zones} di Kabupaten Kebumen.

\begin{table}[H]
\centering
\caption{Hasil Rekapitulasi Kinerja 13 Model Downscaling Curah Hujan (Protokol Spatial Block CV dengan Buffer Zone $5\text{ km}$).}
\label{tab:benchmark}
\resizebox{\textwidth}{!}{%
\begin{tabular}{lrrrrrr}
\toprule
\textbf{Algoritma / Model} & \textbf{KGE Mean} & \textbf{RMSE (mm)} & \textbf{MAE (mm)} & \textbf{PBIAS (\%)} & \textbf{OOF-KGE} & \textbf{OOF-RMSE} \\
\midrule
\rowcolor{green!10} \textbf{Ensemble Mean} & \textbf{-0.1019} & 21.13 & 19.03 & +5.04 & 0.3384 & 31.67 \\
\rowcolor{green!5}  \textbf{PRISM} & -0.1297 & \textbf{12.78} & 10.70 & \textbf{-0.60} & \textbf{0.7189} & \textbf{14.18} \\
Spatial XGBoost & -0.1829 & 12.67 & \textbf{9.42} & -1.06 & 0.6037 & 14.41 \\
Regression-Kriging (RK) & -0.1984 & 18.93 & 15.14 & +2.12 & 0.3987 & 24.54 \\
KED (Elevation Drift) & -0.2397 & 13.36 & 10.60 & +0.14 & 0.6603 & 15.21 \\
GWR (Elevation Local) & -0.2432 & 12.31 & 9.64 & +0.16 & 0.7305 & 13.72 \\
Spatial Random Forest & -0.2842 & 14.55 & 10.74 & -1.28 & 0.3063 & 17.07 \\
RFSI & -0.3172 & 16.84 & 13.10 & +1.52 & 0.1868 & 18.19 \\
IDW (Power 2) & -0.3679 & 21.69 & 18.53 & +0.80 & -0.6236 & 23.22 \\
Spatial LightGBM & -0.4163 & 23.00 & 19.95 & +1.59 & -0.8449 & 23.67 \\
ANUSPLIN (Trivariate) & -0.9630 & 62.43 & 55.48 & +19.71 & -2.0322 & 104.31 \\
Ordinary Kriging & -2.1003 & 60.33 & 52.37 & -11.72 & -2.7543 & 109.63 \\
Conditional Merging & -3.1475 & 90.19 & 75.87 & -1.33 & -4.2646 & 128.11 \\
\bottomrule
\end{tabular}%
}
\end{table}

\begin{figure}[H]
    \centering
    \includegraphics[width=0.95\textwidth]{figures/benchmark_13models_evaluation.png}
    \caption{Perbandingan Efisiensi Model (KGE) dan Tingkat Kesalahan (RMSE) pada 13 Algoritma Downscaling Curah Hujan di Bawah Pengujian Spatial Block Buffer.}
    \label{fig:benchmark_bars}
\end{figure}

\subsection{Diskusi Hasil Benchmarking}
Berdasarkan Gambar \ref{fig:benchmark_bars} dan Tabel \ref{tab:benchmark}, beberapa temuan kunci dapat diidentifikasi:
\begin{enumerate}
    \item \textbf{Keunggulan Pendekatan Regresi Berbasis Fisika Orografis (PRISM \& Ensemble Mean):}
    PRISM membuktikan keunggulannya dengan nilai kesalahan absolut terendah ($\text{RMSE} = 12,78\text{ mm}$, $\text{PBIAS} = -0,60\%$, dan $\text{OOF-KGE} = 0,7189$). Hal ini disebabkan oleh mekanisme bobot faset tiga dimensi yang secara ketat mencegah stasiun di dataran pantai membiaskan estimasi di lereng pegunungan. Model Ensemble Mean memuncaki skor rata-rata KGE lintas fold ($\text{KGE} = -0,1019$) karena kemampuannya meredam osilasi liar dari model individual.
    \item \textbf{Kelemahan Geostatistik Murni Tanpa Drift (Ordinary Kriging):}
    Ordinary Kriging mengalami kegagalan signifikan ($\text{KGE} = -2,1003$; $\text{RMSE} = 60,33\text{ mm}$) saat diuji pada blok spasial terisolasi. Karena hanya mengandalkan kovariansi jarak horizontal tanpa mempertimbangkan elevasi medan, OK tidak mampu mengantisipasi lonjakan presipitasi di puncak pegunungan utara ketika stasiun latih berada di dataran rendah.
    \item \textbf{Perilaku Algoritma Machine Learning (XGBoost vs LightGBM vs RF):}
    Spatial XGBoost menunjukkan performa sangat kompetitif ($\text{RMSE} = 12,67\text{ mm}$; $\text{OOF-KGE} = 0,6037$) berkat regularisasi gradien L2 yang efektif. Sebaliknya, model pohon keputusan yang terlalu agresif rentan terhadap \textit{spatial boundary overfitting} jika kedalaman daun tidak dibatasi secara ketat.
\end{enumerate}

% ==============================================================================
\section{ANALISIS SPASIAL DAN TEMPORAL DOWNSCALING TAHUN 2025}
% ==============================================================================

\subsection{Dinamika Temporal dan Trajektori Presipitasi Bulanan 2025 Lintas Algoritma}
Berdasarkan eksekusi downscaling untuk seluruh 12 bulan pada tahun 2025, trajektori temporal presipitasi di Kabupaten Kebumen memperlihatkan konsistensi hidrologis yang sangat tinggi antara data satelit makro (CHIRPS raw $5,5\text{ km}$) dan keempat algoritma downscaling (Tabel \ref{tab:monthly_all_algorithms} dan Gambar \ref{fig:monthly_trajectory}).

\begin{figure}[H]
    \centering
    \includegraphics[width=0.96\textwidth]{figures/monthly_trajectory_algorithms_comparison_2025.png}
    \caption{Perbandingan Trajektori Rata-Rata Curah Hujan Bulanan Lintas Algoritma di Kabupaten Kebumen Tahun 2025. Terlihat pelestarian konservasi massa regional yang sempurna di mana seluruh algoritma bergerak harmonis mengikuti dinamika monsun tropis.}
    \label{fig:monthly_trajectory}
\end{figure}

\begin{table}[H]
\centering
\caption{Perbandingan Rinci Curah Hujan Bulanan Lintas Algoritma di Kabupaten Kebumen Tahun 2025 (Nilai Rata-Rata Wilayah, Rentang Spasial Min--Max, dan Ketidakpastian Inter-Model $\sigma$ dalam satuan mm).}
\label{tab:monthly_all_algorithms}
\resizebox{\textwidth}{!}{%
\begin{tabular}{clrrrrrrr}
\toprule
\textbf{No} & \textbf{Bulan} & \textbf{CHIRPS Raw} & \textbf{ANUSPLIN} & \textbf{PRISM} & \textbf{RegKriging} & \textbf{Ensemble} & \textbf{Rentang Spasial (Min--Max)} & \textbf{Spread $\sigma$} \\
\midrule
01 & Januari   & 430.99 & 422.87 & 424.30 & 423.21 & \textbf{423.46} & 356.77 -- 548.37 & $\pm 5.81$ \\
02 & Februari  & 354.29 & 348.48 & 349.01 & 348.19 & \textbf{348.56} & 265.53 -- 482.15 & $\pm 5.15$ \\
03 & Maret     & 424.25 & 410.12 & 411.08 & 410.82 & \textbf{410.68} & 294.48 -- 686.02 & $\pm 8.94$ \\
04 & April     & 141.43 & 134.21 & 134.16 & 134.47 & \textbf{134.28} &  73.89 -- 266.27 & $\pm 3.72$ \\
05 & Mei       & 448.86 & 448.07 & 446.48 & 449.81 & \textbf{448.12} & 356.81 -- 588.46 & $\pm 7.78$ \\
06 & Juni      & 107.89 & 103.40 & 103.57 & 103.13 & \textbf{103.36} &  37.24 -- 149.94 & $\pm 2.34$ \\
07 & Juli      & 146.30 & 143.99 & 143.71 & 143.55 & \textbf{143.75} &  85.70 -- 211.56 & $\pm 5.30$ \\
08 & Agustus   & 392.80 & 380.66 & 379.99 & 380.69 & \textbf{380.45} & 181.17 -- 522.23 & $\pm 7.68$ \\
09 & September & 164.68 & 162.19 & 161.06 & 162.91 & \textbf{162.05} &  70.56 -- 234.04 & $\pm 4.35$ \\
10 & Oktober   & 375.61 & 379.32 & 376.50 & 381.20 & \textbf{379.01} & 242.78 -- 499.16 & $\pm 7.83$ \\
11 & November  & 316.54 & 323.45 & 320.26 & 324.59 & \textbf{322.77} & 244.82 -- 416.16 & $\pm 6.25$ \\
12 & Desember  & 231.07 & 221.26 & 222.39 & 221.79 & \textbf{221.81} & 143.90 -- 374.76 & $\pm 4.27$ \\
\midrule
\textbf{--} & \textbf{Akumulasi 2025} & \textbf{3534.71} & \textbf{3478.02} & \textbf{3472.51} & \textbf{3484.36} & \textbf{3478.29} & \textbf{2353.65 -- 4979.12} & \textbf{$\pm 69.42$} \\
\bottomrule
\end{tabular}%
}
\end{table}

Puncak presipitasi tahun 2025 terjadi pada bulan Mei ($448,12\text{ mm}$), Januari ($423,46\text{ mm}$), dan Maret ($410,68\text{ mm}$), sementara periode kemarau relatif berlangsung pada bulan Juni ($103,36\text{ mm}$) dan April ($134,28\text{ mm}$). Bulan Agustus mengalami anomali basah substansial ($380,45\text{ mm}$) akibat perambatan konvektif ekuatorial. Yang terpenting, deviasi antar-model ($\sigma$) pada rata-rata regional berada di bawah $9\text{ mm/bulan}$, membuktikan kestabilan numerik luar biasa antar-metode.

\subsection{Pola Spasial 12 Bulan Tiap Algoritma (Analisis Grid 3x4 Masing-Masing Metode)}
Untuk memberikan transparansi penuh mengenai karakteristik perilaku masing-masing algoritma, Gambar \ref{fig:grid_anusplin} hingga \ref{fig:grid_ensemble} menampilkan hasil rekonstruksi spasial 12 bulan resolusi 250m untuk masing-masing model secara terpisah.

\subsubsection{1. Rekonstruksi 12 Bulan Menggunakan ANUSPLIN (Trivariate Topographic Spline)}
Gambar \ref{fig:grid_anusplin} menampilkan rekonstruksi 12 bulan menggunakan ANUSPLIN. Model ini menunjukkan keunggulan khas dalam menghasilkan permukaan isohyet yang sangat kontinu, kelicinan kurvatur yang sempurna, dan pelenyapan total terhadap efek batas piksel satelit, baik pada bulan basah ekstrem maupun bulan kering.

\begin{figure}[H]
    \centering
    \includegraphics[width=0.98\textwidth]{figures/monthly_12m_ANUSPLIN_2025.png}
    \caption{Distribusi Spasial Curah Hujan Bulanan (250m) Tahun 2025 Menggunakan Algoritma ANUSPLIN (Trivariate Topographic Spline).}
    \label{fig:grid_anusplin}
\end{figure}

\subsubsection{2. Rekonstruksi 12 Bulan Menggunakan PRISM (Facet Weighted Orographic Regression)}
Gambar \ref{fig:grid_prism} memperlihatkan rekonstruksi 12 bulan menggunakan PRISM. Berbeda dengan ANUSPLIN yang murni mengandalkan energi kelicinan splines, PRISM secara aktif menyesuaikan bobot lereng faset lokal ($W_d, W_z, W_c$). Struktur kontur isohyet PRISM sangat terikat pada punggungan topografi dan jurang lembah, menghasilkan detail geomorfologi lokal yang sangat tajam di Pegunungan Karangsambung dan Sadang.

\begin{figure}[H]
    \centering
    \includegraphics[width=0.98\textwidth]{figures/monthly_12m_PRISM_2025.png}
    \caption{Distribusi Spasial Curah Hujan Bulanan (250m) Tahun 2025 Menggunakan Algoritma PRISM (Topographic-Coastal Facet Regression).}
    \label{fig:grid_prism}
\end{figure}

\subsubsection{3. Rekonstruksi 12 Bulan Menggunakan Regression-Kriging (RK)}
Gambar \ref{fig:grid_rk} menyajikan hasil 12 bulan menggunakan Regression-Kriging. Terlihat transisi linier yang tegas antara elevasi medan dan residu Gaussian spasial. Pada bulan-bulan transisi (seperti April dan September), RK secara konsisten mempertahankan variogram spasial stasioner lokal.

\begin{figure}[H]
    \centering
    \includegraphics[width=0.98\textwidth]{figures/monthly_12m_RegressionKriging_2025.png}
    \caption{Distribusi Spasial Curah Hujan Bulanan (250m) Tahun 2025 Menggunakan Algoritma Regression-Kriging (Topographic Trend + Gaussian Kriging).}
    \label{fig:grid_rk}
\end{figure}

\subsubsection{4. Rekonstruksi 12 Bulan Konsensus Multi-Model (Ensemble Mean)}
Gambar \ref{fig:grid_ensemble} menampilkan hasil konsensus Ensemble Mean. Model ini berhasil memadukan kelicinan permukaan ANUSPLIN, ketajaman faset orografis PRISM, dan kekuatan kovariansi spasial Regression-Kriging, menciptakan estimasi presipitasi paling seimbang dengan ketidakpastian terendah.

\begin{figure}[H]
    \centering
    \includegraphics[width=0.98\textwidth]{figures/monthly_12m_EnsembleMean_2025.png}
    \caption{Distribusi Spasial Curah Hujan Bulanan (250m) Tahun 2025 Menggunakan Konsensus Ensemble Multi-Model (Ensemble Mean).}
    \label{fig:grid_ensemble}
\end{figure}

\subsection{Matriks Komparasi Spasial 4 Algoritma pada Periode Musim Representatif}
Untuk membandingkan respon keempat model secara berdampingan pada kondisi iklim yang berbeda, Gambar \ref{fig:seasonal_matrix} menyajikan matriks komparasi spasial berukuran $4 \times 4$ yang membedah 4 periode musim representatif:
\begin{enumerate}
    \item \textbf{Januari (Puncak Musim Penghujan Awal):} Ditandai oleh intensitas presipitasi masif di atas $420\text{ mm}$, di mana seluruh model secara serentak mengidentifikasi pusat hujan maksimum di lereng utara.
    \item \textbf{Mei (Puncak Presipitasi Tertinggi):} Akumulasi bulanan mencapai $> 448\text{ mm}$, PRISM dan RK memperlihatkan konsentrasi hujan orografis yang sangat intens di perbukitan Ka\-rang\-ga\-yam.
    \item \textbf{Juli (Periode Musim Kemarau):} Hujan berkurang drastis ($143\text{ mm}$), ANUSPLIN menghasilkan transisi paling mulus tanpa artefak batas kering, sedangkan PRISM memetakan kantong-kantong hujan lokal pada igir pegunungan tertinggi.
    \item \textbf{Oktober (Awal Musim Hujan Transisi):} Presipitasi melonjak kembali ke $379\text{ mm}$, memperlihatkan respon cepat seluruh algoritma terhadap pergeseran massa udara muson barat.
\end{enumerate}

\begin{figure}[H]
    \centering
    \includegraphics[width=0.98\textwidth]{figures/seasonal_matrix_4algorithms_comparison_2025.png}
    \caption{Matriks Komparasi Spasial $4 \times 4$ Lintas 4 Algoritma pada 4 Periode Musim Representatif (Januari, Mei, Juli, dan Oktober 2025) di Kabupaten Kebumen.}
    \label{fig:seasonal_matrix}
\end{figure}

\subsection{Komparasi Spasial 4 Model Terbaik pada Akumulasi Tahunan 2025}
Gambar \ref{fig:comparison_4p} memperlihatkan perbandingan spasial 4 model terbaik dunia (ANUSPLIN, PRISM, Regression-Kriging, dan Ensemble Mean) pada akumulasi curah hujan tahunan 2025.

\begin{figure}[H]
    \centering
    \includegraphics[width=0.98\textwidth]{figures/spatial_4models_comparison_2025.png}
    \caption{Perbandingan Spasial 4 Algoritma Downscaling pada Akumulasi Curah Hujan Tahunan 2025 di Kabupaten Kebumen (Resolusi 250 meter).}
    \label{fig:comparison_4p}
\end{figure}

Analisis komparatif menunjukkan:
\begin{enumerate}
    \item \textbf{ANUSPLIN:} Menghasilkan transisi orografis paling halus (\textit{maximally smooth surface}) dengan gradien kurvatur kontinu, ideal untuk pemodelan visual dan hidrologi regional.
    \item \textbf{PRISM:} Sangat sensitif terhadap faset kelerengan lereng lokal (\textit{facet-driven}), memisahkan secara tegas antara lereng yang menghadap laut (\textit{windward}) dan lereng bayangan hujan (\textit{leeward}).
    \item \textbf{Regression-Kriging:} Menangkap efek tren makro regional secara linier dengan elevasi, sementara variabilitas lokal dikendalikan oleh semivariogram residu.
    \item \textbf{Ensemble Mean Consensus:} Menghasilkan integrasi paling stabil dan seimbang, meminimalkan risiko overestimasi lokal yang mungkin terjadi pada model tunggal.
\end{enumerate}

\subsection{Analisis Ketidakpastian Inter-Model (Uncertainty Spread)}
Keandalan estimasi spasial dianalisis melalui deviasi standar antar-model ($\sigma_{\text{model}}$), sebagaimana ditampilkan pada Gambar \ref{fig:uncertainty_map}.

\begin{figure}[H]
    \centering
    \includegraphics[width=0.92\textwidth]{figures/model_uncertainty_spread_2025.png}
    \caption{Peta Spasial Ketidakpastian Inter-Model ($\sigma$) Akumulasi Tahunan 2025. Zona merah tua menandakan wilayah dengan tingkat dispersi antar-model yang dipicu oleh kompleksitas topografi curam.}
    \label{fig:uncertainty_map}
\end{figure}

Peta ketidakpastian mengungkapkan temuan metodologis penting:
\begin{itemize}
    \item \textbf{Konsensus Sangat Tinggi di Dataran Aluvial Pesisir:} Di seluruh wilayah selatan (Ambal, Mirit, Klirong, Buluspesantren), nilai $\sigma < 35\text{ mm}$ (kurang dari $1,2\%$ dari total hujan tahunan). Ini membuktikan bahwa seluruh model sepakat penuh ketika beroperasi di wilayah dengan variasi topografi minimal.
    \item \textbf{Dispersi Terpusat pada Lereng Curam Pegunungan Utara:} Di Kecamatan Sadang dan Karangsambung, nilai $\sigma$ meningkat hingga mencapai $128\text{ mm}$. Hal ini disebabkan oleh perbedaan formulasi matematis: PRISM menggunakan pembobotan eksponensial elevasi, ANUSPLIN menerapkan penalti kelicinan splines, dan RK mengandalkan variogram Gaussian.
\end{itemize}

\subsection{Statistik Zonal Presipitasi 26 Kecamatan di Kabupaten Kebumen}
Tabel \ref{tab:zonal_kecamatan} menyajikan hasil ekstraksi statistik zonal presipitasi hasil downscaling untuk ke-26 kecamatan di Kabupaten Kebumen sepanjang tahun 2025.

\begin{table}[H]
\centering
\caption{Statistik Zonal Curah Hujan Hasil Downscaling Tahun 2025 di 26 Kecamatan Kabupaten Kebumen.}
\label{tab:zonal_kecamatan}
\resizebox{\textwidth}{!}{%
\begin{tabular}{clrrrrrrr}
\toprule
\textbf{No} & \textbf{Kecamatan} & \textbf{Elevasi (m)} & \textbf{Jarak Pantai (km)} & \textbf{Tahunan (mm)} & \textbf{ANUSPLIN} & \textbf{PRISM} & \textbf{RegKriging} & \textbf{Spread $\sigma$} \\
\midrule
1 & Ayah & 150.2 & 15.0 & 3138.0 & 3121.5 & 3113.7 & 3178.7 & 77.4 \\
2 & Buayan & 106.4 & 18.6 & 3147.5 & 3157.0 & 3215.3 & 3070.2 & 89.8 \\
3 & Puring & 10.1 & 12.6 & 3066.6 & 3061.3 & 3113.4 & 3025.2 & 52.4 \\
4 & Petanahan & 12.4 & 12.5 & 3188.5 & 3177.8 & 3169.4 & 3218.1 & 34.6 \\
5 & Klirong & 15.5 & 12.2 & 3173.9 & 3176.3 & 3165.8 & 3179.7 & 36.1 \\
6 & Buluspesantren & 13.0 & 10.4 & 3033.3 & 3027.6 & 3054.7 & 3017.5 & 39.5 \\
7 & Ambal & 11.7 & 8.8 & 2857.5 & 2828.1 & 2912.5 & 2831.9 & 54.9 \\
8 & Mirit & 10.8 & 7.1 & \textbf{2699.1} & 2662.6 & 2779.3 & 2655.6 & 72.0 \\
9 & Prembun & 22.5 & 16.1 & 3048.9 & 3022.5 & 3108.6 & 3015.5 & 56.5 \\
10 & Kutowinangun & 30.2 & 16.0 & 3165.4 & 3137.8 & 3182.4 & 3175.9 & 39.5 \\
11 & Alian & 103.8 & 24.0 & 3779.2 & 3801.6 & 3738.2 & 3797.8 & 56.3 \\
12 & Kebumen & 25.9 & 18.2 & 3413.7 & 3426.9 & 3391.2 & 3423.1 & 33.2 \\
13 & Pejagoan & 109.3 & 24.3 & 3801.2 & 3833.5 & 3748.1 & 3822.0 & 65.8 \\
14 & Sruweng & 85.6 & 22.6 & 3660.0 & 3671.5 & 3611.9 & 3696.5 & 63.6 \\
15 & Adimulyo & 10.9 & 19.2 & 3365.8 & 3341.7 & 3347.4 & 3408.4 & 50.0 \\
16 & Kuwarasan & 12.3 & 20.8 & 3221.7 & 3238.7 & 3313.8 & 3112.6 & 97.3 \\
17 & Rowokele & 183.8 & 27.3 & 3465.1 & 3456.8 & 3479.7 & 3458.8 & 77.0 \\
18 & Sempor & 170.7 & 33.0 & 3720.4 & 3713.2 & 3744.2 & 3704.0 & 63.7 \\
19 & Gombong & 23.2 & 26.1 & 3475.6 & 3458.0 & 3508.1 & 3460.8 & 47.7 \\
20 & Karanganyar & 82.1 & 26.2 & 3697.2 & 3676.2 & 3653.2 & 3762.2 & 68.4 \\
21 & Karanggayam & 181.8 & 33.2 & 4175.8 & 4198.1 & 4087.2 & 4242.0 & 103.4 \\
22 & Sadang & 291.0 & 38.1 & \textbf{4547.6} & 4593.9 & 4447.7 & 4601.4 & 128.1 \\
23 & Bonorowo & 8.9 & 9.8 & 2723.4 & 2695.5 & 2818.1 & 2656.6 & 74.0 \\
24 & Padureso & 139.9 & 23.8 & 3585.9 & 3589.1 & 3631.9 & 3536.7 & 62.8 \\
25 & Poncowarno & 116.3 & 20.9 & 3493.8 & 3475.6 & 3496.0 & 3510.0 & 46.5 \\
26 & Karangsambung & 178.1 & 30.5 & 4266.2 & 4276.9 & 4142.0 & 4379.6 & 119.5 \\
\bottomrule
\end{tabular}%
}
\end{table}

\begin{figure}[H]
    \centering
    \includegraphics[width=0.95\textwidth]{figures/zonal_kecamatan_hyetograph_2025.png}
    \caption{Profil Presipitasi Bulanan 2025: Bukti Nyata Penguatan Orografis pada Kecamatan Pegunungan Utara (Garis Segitiga) Dibandingkan Kecamatan Pesisir Selatan (Garis Lingkaran Putus-putus).}
    \label{fig:hyetograph_zon}
\end{figure}

Gambar \ref{fig:hyetograph_zon} membuktikan secara empiris fenomena orografis yang selama ini tertutup oleh resolusi kasar satelit:
\begin{itemize}
    \item Tiga kecamatan di zona pegunungan utara mencatat curah hujan tertinggi: \textbf{Sadang} ($4.547,6\text{ mm}$), \textbf{Karangsambung} ($4.266,2\text{ mm}$), dan \textbf{Karanggayam} ($4.175,8\text{ mm}$).
    \item Sebaliknya, kecamatan di pesisir pantai selatan menerima curah hujan paling rendah: \textbf{Mirit} ($2.699,1\text{ mm}$), \textbf{Bonorowo} ($2.723,4\text{ mm}$), dan \textbf{Ambal} ($2.857,5\text{ mm}$).
    \item Selisih presipitasi mencapai $1.848,5\text{ mm/tahun}$ ($+68,5\%$), mengindikasikan bahwa DAS Luk Ulo hulu menerima pasokan air limpasan yang jauh lebih masif daripada yang terdeteksi oleh produk satelit CHIRPS mentah.
\end{itemize}

% ==============================================================================
\section{DISKUSI DAN REKOMENDASI TEKNIS}
% ==============================================================================

\subsection{Keseimbangan Antara Smoothing dan Detail Fisik}
Salah satu tantangan mendasar dalam downscaling presipitasi adalah menghindari jebakan \textit{"bulls-eye artifacts"} (lingkaran konsentris di sekitar stasiun) pada metode geostatistik murni, serta menghindari \textit{"checkerboard boundary jumps"} pada metode interpolasi terdekat. ANUSPLIN dan PRISM terbukti memberikan keseimbangan optimal:
\begin{itemize}
    \item ANUSPLIN menghasilkan permukaan isohyet yang sangat kontinu dan estetis secara visual tanpa sudut patahan tajam.
    \item PRISM mengaitkan setiap perubahan kontur dengan faset lereng aktual, memberikan akurasi fisik tertinggi pada lereng berhadapan langsung dengan angin basah.
    \item Penggabungan keduanya dalam \textbf{Ensemble Mean Consensus} terbukti melenyapkan kelemahan masing-masing algoritma tunggal.
\end{itemize}

\subsection{Konservasi Massa Presipitasi (Mass Conservation)}
Dalam penerapan operasional hidrologi, algoritma downscaling wajib mempertahankan konservasi massa volume air regional. Skema kalibrasi proporsional yang diterapkan dalam pipeline ini:
\begin{equation}
    P_{\text{calibrated}}(x, y) = P_{\text{pred}}(x, y) \times \left( \frac{\bar{P}_{\text{CHIRPS\_domain}}}{\bar{P}_{\text{pred\_domain}}} \right)
\end{equation}
menjamin bahwa total volume air hujan tahunan yang jatuh di atas Kabupaten Kebumen tetap setara dengan estimasi satelit skala luas, namun didistribusikan ulang secara realistis ke lereng-lereng pegunungan.

\subsection{Rekomendasi Implementasi Operasional}
Berdasarkan temuan studi ini, rekomendasi operasional untuk instansi pengelola bencana dan sumber daya air di Kabupaten Kebumen adalah sebagai berikut:
\begin{enumerate}
    \item \textbf{Untuk Pemodelan Hidrologi \& Debit Banjir DAS Luk Ulo:} Disarankan menggunakan raster hasil \textbf{PRISM} atau \textbf{Ensemble Mean}, karena akurasi presipitasi di wilayah hulu terbukti paling realistis.
    \item \textbf{Untuk Pembuatan Peta Atlas Iklim dan Tata Ruang Bappeda:} Disarankan menggunakan raster hasil \textbf{ANUSPLIN}, karena memberikan kehalusan isohyet terbaik tanpa artefak diskontinuitas visual.
    \item \textbf{Untuk Analisis Kerentanan Bencana Longsor BPBD:} Peta ketidakpastian inter-model ($\sigma$) harus digunakan sebagai lapisan pembobot risiko: zona dengan curah hujan tinggi dan ketidakpastian tinggi (seperti tebing terjal Karanggayam dan Sadang) harus diprioritaskan untuk pemasangan sensor \textit{extensometer} pergerakan tanah.
\end{enumerate}

% ==============================================================================
\section{KESIMPULAN}
% ==============================================================================
Penelitian ini telah berhasil mengeksekusi pipeline \textit{spatial downscaling} komprehensif untuk data satelit CHIRPS resolusi tinggi (250 meter) sepanjang tahun 2025 di Kabupaten Kebumen. Dari benchmarking 13 algoritma dengan protokol ketat \textit{Spatial Block Cross-Validation with Buffer Zones} ($5\text{ km}$), model \textbf{Ensemble Mean Consensus} dan \textbf{PRISM} terbukti sebagai algoritma paling tangguh, unggul, dan konsisten secara hidrologis. 

Peta hasil downscaling tahun 2025 membuktikan kemampuannya dalam melenyapkan artefak piksel kasar sekaligus mengungkap pola orografis mikro yang intensif: wilayah pegunungan utara (Sadang) menerima akumulasi curah hujan tahunan mencapai $4.547,6\text{ mm}$, berbanding kontras dengan wilayah pesisir selatan (Mirit) yang hanya menerima $2.699,1\text{ mm}$. Seluruh 17 raster GeoTIFF beresolusi 250m, tabel statistik zonal 26 kecamatan, dan visualisasi beresolusi cetak 300 DPI yang dihasilkan siap digunakan untuk mendukung tata kelola sumber daya air dan mitigasi bencana di Kabupaten Kebumen.

% ==============================================================================
% DAFTAR PUSTAKA
% ==============================================================================
\begin{thebibliography}{99}

\bibitem{funk2015chirps}
Funk, C., Peterson, P., Landsfeld, M., Pedreros, D., Verdin, J., Shukla, S., Husak, G., Rowland, J., Harrison, L., Hoell, A., \& Michaelsen, J. (2015).
The climate hazards infrared precipitation with stations—a new environmental library for monitoring extremes.
\textit{Scientific Data}, 2(1), 150066.

\bibitem{hutchinson1995anusplin}
Hutchinson, M. F. (1995).
Interpolating mean rainfall using thin plate smoothing splines.
\textit{International Journal of Geographical Information Systems}, 9(4), 385--403.

\bibitem{daly2008physiographically}
Daly, C., Halbleib, M., Smith, J. I., Gibson, W. P., Doggett, M. K., Taylor, G. H., Curtis, J., \& Pasteris, P. P. (2008).
Physiographically sensitive mapping of climatological temperature and precipitation across the conterminous United States.
\textit{International Journal of Climatology}, 28(15), 2031--2064.

\bibitem{hengl2018random}
Hengl, T., Nussbaum, M., Wright, M. N., Heuvelink, G. B., \& Gr{\"a}ler, B. (2018).
Random forest as a generic framework for predictive pseudo-linear and non-linear geostatistical mapping.
\textit{PeerJ}, 6, e5518.

\bibitem{gupta2009decomposition}
Gupta, H. V., Kling, H., Yilmaz, K. K., \& Martinez, G. F. (2009).
Decomposition of the mean squared error and NSE performance criteria: Implications for improving hydrological modelling.
\textit{Journal of Hydrology}, 377(1-2), 80--91.

\bibitem{roberts2017cross}
Roberts, D. R., Bahn, V., Ciuti, S., Boyce, M. S., Elith, J., Guillera-Arroita, G., Hauenstein, S., Lahoz-Monfort, J. J., Schr{\"o}der, B., Thuiller, W., \& Warton, D. I. (2017).
Cross-validation strategies for data with spatial, temporal, hierarchical or phylogenetic structure.
\textit{Ecography}, 40(8), 913--929.

\bibitem{sinclair2005conditional}
Sinclair, S., \& Pegram, G. (2005).
Combining radar and rain gauge estimates of rainfall using conditional merging.
\textit{Atmospheric Science Letters}, 6(1), 19--22.

\bibitem{tobler1970first}
Tobler, W. R. (1970).
A computer movie simulating urban growth in the Detroit region.
\textit{Economic Geography}, 46(sup1), 234--240.

\end{thebibliography}

\end{document}
